\documentclass[11pt]{article}
\usepackage[a4paper,margin=27mm]{geometry}
\usepackage{amsmath,amssymb,mathtools}
\usepackage{fontspec,unicode-math}
\setmainfont{Latin Modern Roman}
\setmathfont{Latin Modern Math}
\usepackage{xurl}
\usepackage[hidelinks]{hyperref}
\newcommand{\sref}[2]{\hyperlink{source:#1:#2}{[S#2]}}
\newcommand{\cataloguescope}[1]{\paragraph{Mathematical question.}#1}
\hypersetup{pdftitle={Problem statement}}
\begin{document}
\section{Wikipedia number-theory item 109: The uniqueness conjecture for Markov numbers that every Markov number is the largest number in exact…}\label{statement}

\subsection{Definitions and mathematical statement}

\paragraph{Shared notation.}$\mathbb N=\{1,2,\ldots\}$, and $\mathbb Z,\mathbb Q,\mathbb R,\mathbb C$ denote the integers, rationals, reals and complex numbers. Any other conventions are specified locally.


A normalized Markov triple is $(a,b,c)\in\mathbb N^3$ with $a\le b\le c$ and $a^2+b^2+c^2=3abc$. A Markov number is any component of a positive integer solution of this equation; the mutation $(a,b,c)\mapsto(a,b,3ab-c)$ and permutations show that every Markov number occurs as a largest component. Normalization removes permutations, not triples with equal entries.
\paragraph{Statement recorded in the supplied source.}
For every Markov number $c$, is there exactly one pair $(a,b)$ with $1\le a\le b\le c$ and $a^2+b^2+c^2=3abc$? \sref{statement}{2}
\subsection{Short English statement}
Does the largest entry uniquely determine a positive Markov triple once the entries are sorted?
\subsection{Sources}
\begin{description}
\item[\hypertarget{source:statement:1}{S1}] ULAM.AI and the UnsolvedMath Contributors, UnsolvedMath: A Curated Collection of Open Mathematics Problems, catalogue entry. CC BY 4.0. \url{https://huggingface.co/datasets/ulamai/UnsolvedMath}.
\item[\hypertarget{source:statement:2}{S2}] Yasuaki Gyoda, \emph{An Introduction to the Lagrange and Markov Spectra through the Lens of Generalized Markov Numbers}, arXiv:2609.28097v1 (2026) Chapter 1, Markov equation and Conjecture 1.1.20. \url{https://arxiv.org/pdf/2609.28097}.
\item[\hypertarget{source:statement:3}{S3}] Wikipedia contributors, \emph{List of unsolved problems in mathematics}, revision 1366636767 (2026) Number theory, Wikipedia number-theory item 109: The uniqueness conjecture for Markov numbers that every Markov number is the largest number in exact…. \url{https://en.wikipedia.org/w/index.php?title=List_of_unsolved_problems_in_mathematics&oldid=1366636767#Number_theory}.
\end{description}
\label{end:statement}
\end{document}
